\documentclass[11pt]{article}

% --- Packages ---
\usepackage[margin=1in]{geometry}
\usepackage[version=4]{mhchem}
\usepackage{chemfig}
\usepackage{amsmath, amssymb}
\usepackage{siunitx}
\usepackage{tcolorbox}
\usepackage{enumitem}
\usepackage{titlesec}
\usepackage{booktabs}
\usepackage{pgfplots}
\pgfplotsset{compat=1.18}

% --- Style Configuration ---
\sisetup{per-mode=symbol}

% Custom tcolorbox styles for distinguishing Provided vs Must Memorize
\tcbset{
    base/.style={fonttitle=\bfseries, left=5pt, right=5pt, top=5pt, bottom=5pt},
    provided/.style={base, colback=blue!5!white, colframe=blue!75!black, title=Key Formula (Provided on Exam)},
    memorize/.style={base, colback=gray!5!white, colframe=gray!50!black, title=Key Formula (Must Memorize)}
}

\titleformat{\section}{\large\bfseries\sffamily}{Chapter \thesection:}{1em}{}
\titleformat{\subsection}{\normalsize\bfseries\sffamily}{}{0em}{}

\begin{document}

\title{CHM2045C: Final Exam Comprehensive Study Guide}
\author{Chemistry Study Guide Draft}
\date{Spring 2026}
\maketitle

\section{Foundations: The Chemical Toolbox}

This chapter establishes the fundamental quantitative and qualitative tools required for chemical analysis, focusing on measurement, classification, and naming.

\subsection{Density and Volume Displacement}
Density is an intrinsic property relating mass and volume. For irregular solids, volume is determined by the displacement of a liquid.
\begin{tcolorbox}[memorize]
    \[ d = \frac{m}{V} \quad \text{and} \quad V_{\text{object}} = V_{\text{final}} - V_{\text{initial}} \]
\end{tcolorbox}
\begin{enumerate}
    \item Identify given mass (\si{g}) and volume (\si{mL} or \si{cm^3}).
    \item Use displacement to find $V$ if necessary (Archimedes' Principle).
    \item Solve for the unknown variable using algebraic rearrangement.
\end{enumerate}

\subsection{Chemical Nomenclature}
\begin{itemize}
    \item \textbf{Ionic Compounds}: Name the cation followed by the anion. Use Roman numerals for transition metals with variable oxidation states (e.g., \ce{FeCl3} is iron(III) chloride).
    \item \textbf{Acids}: Binary acids (\ce{H} + element) use the prefix \textit{hydro-} and suffix \textit{-ic} (e.g., \ce{HCl(aq)} is hydrochloric acid). Oxyacids use \textit{-ic} for \textit{-ate} ions and \textit{-ous} for \textit{-ite} ions.
    \item \textbf{Hydrates}: Append the number of water molecules using Greek prefixes (e.g., \ce{CuSO4.5H2O} is copper(II) sulfate pentahydrate).
\end{itemize}

\subsection{Classification of Matter}
Matter is classified by its composition and uniformity at the particulate level.
\begin{itemize}
    \item \textbf{Elements}: Pure substances consisting of one type of atom (e.g., \ce{O2}).
    \item \textbf{Compounds}: Pure substances with fixed ratios of different atoms (e.g., \ce{H2O}).
    \item \textbf{Mixtures}: Physical blends of substances. \textbf{Homogeneous} (uniform) or \textbf{Heterogeneous} (non-uniform).
\end{itemize}
\begin{center}
    \textit{Particulate Representation of a Water Molecule (\ce{H2O}):} \\
    \vspace{0.5em}
    \chemfig{O(-[1]H)(-[7]H)}
\end{center}

\subsection{Metric System, Units, and Temperature}
Mastery of prefixes is essential for dimensional analysis and unit comparison.
\begin{table}[h]
    \centering
    \begin{tabular}{lcccc}
        \toprule
        Prefix & Nano (n) & Micro ($\mu$) & Milli (m) & Kilo (k) \\
        \midrule
        Factor & $10^{-9}$ & $10^{-6}$ & $10^{-3}$ & $10^{3}$ \\
        \bottomrule
    \end{tabular}
\end{table}

\begin{tcolorbox}[provided, title=Temperature Conversion (Provided)]
    \[ T_{\si{\kelvin}} = T_{^\circ\text{C}} + 273.15 \]
\end{tcolorbox}
\textbf{Note}: Kelvin is mandatory for all gas law and thermodynamic calculations. Absolute zero is \SI{0}{\kelvin}.

\newpage

\section{Atomic Structure}

Understanding the atom involves identifying its constituent subatomic particles and their contribution to the element's identity and mass.

\subsection{Nuclide Notation and Subatomic Particles}
The identity of an atom is defined by its atomic number ($Z$), while its mass is determined by the sum of protons and neutrons.
\begin{tcolorbox}[memorize]
    \[ A = Z + N \quad \text{and} \quad \text{Charge} = \text{Protons} - \text{Electrons} \]
\end{tcolorbox}
In the symbol \ce{^{A}_{Z}X^{charge}}:
\begin{itemize}
    \item \textbf{Protons}: Equal to $Z$ (Atomic Number). Defines the element.
    \item \textbf{Neutrons}: Calculated as $A - Z$ (Mass Number - Atomic Number).
    \item \textbf{Electrons}: Calculated as $Z - (\text{Charge})$. In neutral atoms, $e^- = p^+$.
\end{itemize}

\subsection{Average Atomic Mass}
The mass reported on the periodic table is a weighted average of all naturally occurring isotopes based on their abundance.
\begin{tcolorbox}[memorize]
    \[ \text{Avg. Atomic Mass} = \sum_{i} (\text{Mass of Isotope}_i \times \text{Fractional Abundance}_i) \]
\end{tcolorbox}
\textbf{Technique}: Convert percentages to decimals (e.g., $75\% \rightarrow 0.75$) to find the "fractional abundance" before multiplying by isotopic mass.

\subsection{Isotopes vs. Ions}
\begin{itemize}
    \item \textbf{Isotopes}: Atoms of the same element (same $Z$) with different numbers of neutrons ($N$), resulting in different mass numbers ($A$).
    \item \textbf{Ions}: Atoms that have gained or lost electrons, resulting in a net positive (cation) or negative (anion) charge.
\end{itemize}

\newpage

\section{Electronic Structure and Periodic Trends}

\subsection{Electromagnetic Radiation and the Bohr Model}
Light behaves as both a wave and a particle (photon). The energy of an electron in an atom is quantized.
\begin{tcolorbox}[memorize, title=Light and Photon Energy (Must Memorize)]
    \[ c = \lambda \nu \quad \text{and} \quad E = h\nu = \frac{hc}{\lambda} \]
    \[ \Delta E = -2.18 \times 10^{-18} \si{J} \left( \frac{1}{n_f^2} - \frac{1}{n_i^2} \right) \]
\end{tcolorbox}
Note: $h$ and $c$ are usually provided, but the relationships between $E$, $\lambda$, and $\nu$ must be mastered.

\subsection{Quantum Numbers}
Four quantum numbers describe the unique quantum state (or "address") of an electron:
\begin{enumerate}
    \item \textbf{Principal ($n$)}: Energy level ($1, 2, 3 \dots$).
    \item \textbf{Angular Momentum ($l$)}: Orbital shape ($0$ to $n-1$). $s=0, p=1, d=2, f=3$.
    \item \textbf{Magnetic ($m_l$)}: Orbital orientation ($-l$ to $+l$).
    \item \textbf{Spin ($m_s$)}: Direction of electron spin ($+1/2$ or $-1/2$).
\end{enumerate}

\subsection{Electron Configurations}
Follow three core principles to populate orbitals in the ground state:
\begin{itemize}
    \item \textbf{Aufbau Principle}: Electrons fill the lowest energy orbitals first ($1s, 2s, 2p \dots$).
    \item \textbf{Pauli Exclusion Principle}: No two electrons can have the same four quantum numbers (max 2 $e^-$ per orbital with opposite spins).
    \item \textbf{Hund's Rule}: Electrons fill degenerate orbitals singly with parallel spins before pairing.
\end{itemize}
\textbf{Noble Gas Shorthand}: Use the preceding noble gas to represent core electrons (e.g., \ce{Sn} is \ce{[Kr] 5s^2 4d^{10} 5p^2}).

\subsection{Periodic Trends: Ionization Energy (IE)}
Ionization Energy is the energy required to remove an electron from a gaseous atom.
\begin{itemize}
    \item \textbf{Across a Period}: IE generally increases due to increasing effective nuclear charge ($Z_{\text{eff}}$).
    \item \textbf{Down a Group}: IE generally decreases due to increasing atomic radius and electronic shielding.
\end{itemize}

\newpage

\section{Formula Calculations and the Mole}

\subsection{Core Concepts}
The mole is the fundamental unit in chemistry for measuring the amount of a substance. It provides a bridge between the microscopic world of atoms and molecules and the macroscopic world of grams and liters.

\begin{itemize}[noitemsep]
    \item \textbf{Avogadro's Number}: $1 \text{ mole} = \num{6.022e23}$ representative particles.
    \item \textbf{Molar Mass}: The mass in grams of one mole of a substance (\si{g.mol^{-1}}).
    \item \textbf{Percent Composition}: The percentage by mass of each element in a compound.
\end{itemize}

\begin{tcolorbox}[provided, title=Avogadro Constant (Provided)]
    \[ N_A = \num{6.022e23} \si{\text{ particles}.mol^{-1}} \]
\end{tcolorbox}

\begin{tcolorbox}[memorize, title=Mole Calculations (Must Memorize)]
    \[ n = \frac{m}{MM} \quad \text{and} \quad N = n \times N_A \]
\end{tcolorbox}

\subsection{Problem-Solving Skills}
\begin{enumerate}
    \item \textbf{Mass-to-Particle Conversions}:
    \begin{itemize}
        \item $\text{Mass (g)} \xrightarrow{\div MM} \text{Moles} \xrightarrow{\times N_A} \text{Particles}$
    \end{itemize}
    
    \item \textbf{Calculating Empirical Formulas}:
    \begin{enumerate}[label=\alph*)]
        \item Assume a \SI{100}{\gram} sample (convert \% directly to grams).
        \item Convert grams of each element to moles using atomic masses.
        \item Divide all mole values by the smallest mole value obtained.
        \item If necessary, multiply by a small integer to reach whole numbers.
    \end{enumerate}

    \item \textbf{Percent Mass Composition}:
    \begin{tcolorbox}[memorize]
        \[ \text{Mass \% of element} = \left( \frac{\text{mass of element in 1 mol compound}}{\text{molar mass of compound}} \right) \times 100\% \]
    \end{tcolorbox}
\end{enumerate}

\newpage

\section{Stoichiometry}

\subsection{Core Concepts}
Stoichiometry involves using the relationships between reactants and products in a balanced chemical equation.

\begin{tcolorbox}[memorize, title=Reaction Efficiency (Must Memorize)]
    \[ \text{Percent Yield} = \left( \frac{\text{Actual Yield}}{\text{Theoretical Yield}} \right) \times 100\% \]
\end{tcolorbox}

\subsection{Problem-Solving Skills}
\begin{enumerate}
    \item \textbf{Mass-to-Mass Stoichiometry}:
    \begin{itemize}
        \item $\text{Mass A} \xrightarrow{\div MM_A} \text{Moles A} \xrightarrow{\text{Mole Ratio}} \text{Moles B} \xrightarrow{\times MM_B} \text{Mass B}$
    \end{itemize}

    \item \textbf{Identifying the Limiting Reactant}:
    \begin{enumerate}[label=\alph*)]
        \item Convert the given amounts of all reactants to moles.
        \item Divide the moles of each reactant by its coefficient in the balanced equation.
        \item The reactant with the smallest resulting value is the \textbf{limiting reactant}.
    \end{enumerate}
\end{enumerate}

\newpage

\section{Solutions and Aqueous Reactions}

\subsection{Core Concepts}
\begin{tcolorbox}[memorize, title=Concentration and Dilution (Must Memorize)]
    \[ M = \frac{n}{V} \quad \text{and} \quad M_1V_1 = M_2V_2 \]
\end{tcolorbox}

\subsection{Solubility Reference}
\begin{table}[h]
    \centering
    \caption{Solubility Guidelines for Ionic Compounds (Must Memorize)}
    \begin{tabular}{ll}
        \toprule
        \textbf{Always Soluble} & \ce{NO3-}, \ce{C2H3O2-}, Group 1 (\ce{Li+}, \ce{Na+}, \ce{K+}), \ce{NH4+} \\
        \midrule
        \textbf{Generally Soluble} & \ce{Cl-}, \ce{Br-}, \ce{I-} (except \ce{Ag+}, \ce{Pb^{2+}}, \ce{Hg2^{2+}}) \\
                                  & \ce{SO4^{2-}} (except \ce{Ba^{2+}}, \ce{Sr^{2+}}, \ce{Ca^{2+}}, \ce{Pb^{2+}}, \ce{Ag+}) \\
        \midrule
        \textbf{Generally Insoluble} & \ce{S^{2-}}, \ce{OH-} (except Group 1, \ce{NH4+}; \ce{Ca^{2+}}, \ce{Sr^{2+}}, \ce{Ba^{2+}} are soluble) \\
                                    & \ce{CO3^{2-}}, \ce{PO4^{3-}} (except Group 1 and \ce{NH4+}) \\
        \bottomrule
    \end{tabular}
\end{table}

\subsection{Oxidation States}
\begin{itemize}
    \item \ce{F} = -1; \ce{O} = -2 (usually); \ce{H} = +1 (with non-metals).
    \item The sum of oxidation states must equal the overall charge.
\end{itemize}

\newpage

\section{Heat and Enthalpy}

\subsection{Calorimetry}
\begin{tcolorbox}[provided, title=Specific Heat Equation (Provided)]
    \[ q = m c_s \Delta T \]
\end{tcolorbox}

\subsection{Reaction Enthalpy Calculations}
\begin{tcolorbox}[memorize, title=Enthalpy Calculations (Must Memorize)]
    \textbf{From Heats of Formation}:
    \[ \Delta H_{\text{rxn}}^\circ = \sum n\Delta H_f^\circ(\text{products}) - \sum m\Delta H_f^\circ(\text{reactants}) \]
    \vspace{0.2em}
    \textbf{From Bond Energies}:
    \[ \Delta H_{\text{rxn}}^\circ = \sum (\text{bonds broken}) - \sum (\text{bonds formed}) \]
\end{tcolorbox}
\textbf{Crucial Distinction}: Formation uses (Products - Reactants), while Bond Energy uses (Reactants/Broken - Products/Formed).

\newpage

\section{Structure and Bonding}

\subsection{Lattice Energy}
Lattice energy increases with increasing ion charge and decreasing ion size.
\begin{tcolorbox}[memorize, title=Coulomb's Law (Must Memorize)]
    \[ E_{\text{lattice}} \propto \frac{Q_1 Q_2}{d} \]
\end{tcolorbox}

\subsection{Formal Charge}
\begin{tcolorbox}[memorize, title=Formal Charge (Must Memorize)]
    \[ FC = (\text{Valence } e^-) - (\text{Non-bonding } e^-) - \frac{1}{2}(\text{Bonding } e^-) \]
\end{tcolorbox}

\newpage

\section{States of Matter}

\subsection{Gas Laws and Constants}
\begin{tcolorbox}[provided, title=Gas Laws and Constants (Provided)]
    \textbf{Constants}:
    \[ R = \SI{0.0821}{\liter.atm.mol^{-1}.K^{-1}} \]
    \[ 1 \text{ atm} = 760 \text{ mmHg} = 760 \text{ Torr} \]
    \vspace{0.5em}
    \textbf{Laws}:
    \[ PV = nRT \quad \text{and} \quad \frac{P_1V_1}{n_1T_1} = \frac{P_2V_2}{n_2T_2} \]
    \vspace{0.5em}
    \textbf{Graham's Law of Effusion}:
    \[ \frac{\text{rate}_1}{\text{rate}_2} = \sqrt{\frac{M_2}{M_1}} \]
\end{tcolorbox}

\subsection{Gas Density}
\begin{tcolorbox}[memorize, title=Gas Density (Must Memorize)]
    \[ d = \frac{PM}{RT} \]
\end{tcolorbox}

\newpage

\section*{Addendum: Exam Formula \& Constant Reference}

\subsection*{Provided Information}
\begin{itemize}
    \item \textbf{Temperature}: $T_{\si{\kelvin}} = T_{^\circ\text{C}} + 273.15$
    \item \textbf{Avogadro Constant}: $N_A = \num{6.022e23} \si{mol^{-1}}$
    \item \textbf{Gas Constant}: $R = \SI{0.0821}{\liter.atm.mol^{-1}.K^{-1}}$
    \item \textbf{Pressure}: $1 \text{ atm} = \SI{760}{mmHg} = \SI{760}{Torr}$
    \item \textbf{Ideal Gas Law}: $PV = nRT$
    \item \textbf{Combined Gas Law}: $\frac{P_1V_1}{n_1T_1} = \frac{P_2V_2}{n_2T_2}$
    \item \textbf{Graham's Law}: $\frac{\text{rate}_1}{\text{rate}_2} = \sqrt{\frac{M_2}{M_1}}$
    \item \textbf{Specific Heat}: $q = m c_s \Delta T$
\end{itemize}

\subsection*{Information to Memorize}
\begin{itemize}
    \item \textbf{Atomic Mass}: $\sum (\text{mass} \times \text{fraction abundance})$
    \item \textbf{Light}: $c = \lambda \nu$ and $E = h\nu = hc/\lambda$
    \item \textbf{Rydberg}: $\Delta E = -2.18 \times 10^{-18} \si{J} (1/n_f^2 - 1/n_i^2)$
    \item \textbf{Molarity}: $M = n/V$
    \item \textbf{Enthalpy (Formation)}: $\sum \Delta H_f^\circ(\text{prod}) - \sum \Delta H_f^\circ(\text{react})$
    \item \textbf{Enthalpy (Bonds)}: $\sum (\text{broken}) - \sum (\text{formed})$
    \item \textbf{Lattice Energy}: $E \propto Q_1Q_2/d$
    \item \textbf{Solubility}: (See Table in Chapter 6)
\end{itemize}

\end{document}
